% =====================================================================
% Chapter 1: Introduction  --  DRAFT v2 (reframed for actual results)
%
% ⚠️ 與 v1 的差別：v1 把 A3（融合選幀）寫成貢獻一。實驗證偽了 A3，
%    因此貢獻清單、RQ、範圍全部重寫。現在的定位是：
%      - proposed method = A1B1（像素選幀 + 雙路決策）
%      - A3 = 被證偽的設計，作為帶機制的負面結果呈現
%
% 標註等級：[OK] / [DERIVED] / [VERIFY] / [FILL]
% =====================================================================

\chapter{Introduction}
\label{chap:introduction}


\section{Background and Motivation}
\label{sec:motivation}

Modern web applications are increasingly defined by their dynamic
behaviour. Interface state is no longer a sequence of discrete page loads
but a continuous flow of transitions: menus that slide, skeleton screens
that resolve into content, components that re-render in response to
asynchronous data. This shift has introduced a class of defect that is
poorly served by existing quality assurance practice. A dropdown whose
transition stutters, a component that flickers during re-rendering, or a
layout that shifts as content loads and causes a user to click the wrong
control are all failures a user experiences directly, yet none is reliably
captured by the tools currently used to test web interfaces.

Two established approaches exist, and each misses this class of defect for
a different reason. Functional testing frameworks such as Playwright and
Selenium assert over the document object model (DOM), and therefore pass
whenever the structural state eventually converges on the expected value,
irrespective of how the interface behaved on the way there. Visual
regression testing (VRT) compares rendered screenshots against approved
baselines, and therefore observes appearance rather than structure, but
does so at isolated instants and with an oracle that cannot distinguish
meaningful change from rendering noise.

The severity of this second limitation is now quantified.
\citet{zhang2026beyond} evaluate pixel-level comparison on a benchmark of
human-verified Web UI changes and find that a naive pixel diff suppresses
none of the non-meaningful changes at all, and that adding a difference
threshold with anti-aliasing tolerance raises suppression only to
\SI{6.72}{\percent} --- a false-positive rate above \SI{93}{\percent}.
The practical consequence is alert fatigue: reviewers confronted with
large volumes of spurious differences at every release cycle come to
distrust the signal.
% [OK] 數字出自 Zhang et al. (2026) Table VI，Claude 已核對。
% [VERIFY] "alert fatigue" 這句目前沒有引用支撐。建議補一筆工業案例研究，
%   或維持這種較保守的措辭。

Vision-language models (VLMs) appear to offer the semantic discrimination
the pixel oracle lacks. A model that can recognise that a button has
shrunk because the viewport narrowed, rather than because a stylesheet
broke, would address the false-positive problem at its source. Yet the
evidence to date suggests that supplying a VLM with raw video and asking
whether a defect is present does not work. TempGlitch reports that VLMs
given temporally sampled video perform close to chance on defect
detection, falling into one of two extremes: flagging almost everything or
almost nothing~\citep{tempglitch}. The same bimodal failure has been
observed in the static web UI setting, where one off-the-shelf model
suppressed \SI{99.63}{\percent} of visual noise but detected only
\SI{65.68}{\percent} of genuine changes, while another failed in the
opposite direction~\citep{zhang2026beyond}.
% [!! RED] ★ TempGlitch 這一筆 Claude 完全沒有查證過 —— 沒有作者、年份、
%   venue、DOI。這是你整篇論文的出發點，必須在提交前親自讀過原文，
%   確認上述結論確實出自該文。詳見 Chapter 2 的說明。
% [OK] Zhang et al. 的兩個數字已核對。

This dissertation proceeds from the hypothesis that the problem lies less
in the model than in what the model is shown and what it is asked. If a
VLM is given a sequence of largely redundant frames sampled without regard
to when anything happened, and asked an unanchored question with no
reference point, poor performance is unsurprising. The intervention
proposed here is therefore not a new model but a new way of constructing
the evidence a model receives and the decision it is asked to make.


\section{Problem Statement}
\label{sec:problem}

The problem addressed by this dissertation may be stated as follows.
Given a recording of a user interacting with a web interface, together
with the structural events emitted by that interface during the
interaction, determine whether the interaction exhibited a dynamic visual
defect, and if so of what kind, at a false positive rate low enough for
the result to be actionable and at a computational cost low enough for the
method to run on every release.

Three properties of this problem make it difficult. First, the evidence is
temporal: a defect may be visible for a fraction of a second within a
recording of several seconds, so the moment of sampling determines whether
it is observed at all. Second, the negative class is not empty of change:
a correctly functioning interface also moves, so an oracle cannot treat
motion as evidence of failure. Third, the cost of inference scales with
the number of frames supplied, so any approach that improves coverage by
supplying more frames trades accuracy against expense.


\section{Research Questions}
\label{sec:rqs}

This dissertation addresses four research questions, each corresponding to
a contrast within the experimental design set out in
Chapter~\ref{chap:methodology}.

\begin{description}
  \item[RQ1] Does adaptive keyframe selection improve detection accuracy
    over uniform temporal sampling, when the reasoning model and the
    prompt are held constant?

  \item[RQ2] Does a dual-path decision procedure, in which the model's
    semantic judgement is combined with structural evidence, reduce the
    false positive rate relative to an unanchored prompt, when the
    keyframe selection strategy is held constant?

  \item[RQ3] Does incorporating DOM mutation events into the keyframe
    selection criterion improve upon selection by pixel difference alone?

  \item[RQ4] What is the resource cost of each condition, measured in
    tokens, latency and monetary terms, relative to the uniform-sampling
    baseline?
\end{description}
% [VERIFY] 這四個 RQ 對應到你的四格：
%   RQ1 = A0B0 vs A3B0（選幀，prompt 固定）
%   RQ2 = A3B0 vs A3B1（決策，選幀固定）
%   RQ3 = A1B1 vs A3B1（融合 vs 純像素，prompt 固定）
%   RQ4 = 全部條件的資源比較
%   請確認這個對應與 Chapter 3、Chapter 4 完全一致。


\section{Contributions}
\label{sec:contributions}

The contributions of this dissertation are as follows. Two are positive
findings and one is a negative finding; all three are reported as
obtained.

\begin{enumerate}
  \item A \textbf{dual-path decision procedure} that treats the model's
    semantic judgement and the interface's structural signals as separate
    lines of evidence and fuses them at the point of decision. This
    procedure eliminated false positives entirely, from 10 of 60 nominal
    clips under the baseline to 0 of 60, in every repetition and under
    both keyframe selectors tested, while raising rather than lowering
    anomaly recall.

  \item A demonstration that \textbf{adaptive keyframe selection is a cost
    intervention rather than an accuracy intervention}. Replacing uniform
    six-frame sampling with three selected frames reduced token
    consumption by between 39 and 49 per cent with no measurable effect on
    accuracy in either direction, a null result of practical value in
    itself.

  \item A \textbf{negative result concerning structural signals in
    selection}. Incorporating DOM mutation events into the keyframe
    selection criterion, which was the design this study originally set
    out to validate, conferred no benefit over pixel-only selection and
    produced a statistically supported degradation when combined with the
    dual-path prompt (all eight discordant clips favouring pixel-only
    selection, $p_{\text{Holm}} = 0.031$, Cohen's $g = 0.500$). The
    finding motivates a distinction between the value of structural
    evidence at the point of decision, where it proved effective, and at
    the point of selection, where it proved harmful.

  \item A \textbf{fully automated fault injection pipeline} that produces
    labelled recordings of dynamic web UI defects without manual
    annotation, together with a quality gate that verifies each injected
    defect is visually manifest before the recording enters the dataset.

  \item An \textbf{empirical evaluation} under a predefined statistical
    protocol, comprising 1{,}440 model invocations across four conditions
    with three repetitions, reporting dispersion, bootstrap interval
    estimates, omnibus and corrected pairwise testing, and effect sizes
    throughout.
\end{enumerate}
% [OK] 第 1、2、3、5 點的數字皆取自 summary.json，Claude 已核對。
% [VERIFY] 第 3 點是本論文最不尋常也最有價值的一點。請確認你願意
%   把「證偽自己原本的設計」明確列為貢獻 —— 這是正確且加分的做法，
%   但你必須在 Demo 面談時能自信地說明為什麼它是貢獻。


\section{Scope and Delimitations}
\label{sec:scope}

Three boundaries are set explicitly at the outset.

The reasoning model is held constant across all experimental conditions. A
single model, \texttt{gemini-3.1-flash-lite}, was used throughout with
temperature fixed at zero. The contribution claimed here concerns the
construction of evidence supplied to a model rather than the relative
capability of models, and holding the model fixed isolates the effect of
interest. Cross-model comparison is identified as future work in
Section~\ref{sec:futurework}.
% [VERIFY] ⚠️ 一致性檢查：你在第二次會議的信裡向教授提過 Qwen2.5-VL，
%   並用它申請了 vLab GPU；你已在 8/2 的信裡主動說明砍掉這一項。
%   請確認 Chapter 3 §3.4 與 Chapter 5 §5.5 的措辭與此處一致，
%   且 Chapter 3 已移除「兩個模型」的舊敘述。

The evaluation is conducted on synthetically generated recordings of a
controlled application served locally. This affords exact ground truth and
reproducible defect timing, both necessary for the paired comparisons
undertaken here, at the cost of external validity. The extent to which the
findings transfer to production interfaces is not established and is
discussed in Section~\ref{sec:limitations}.

The defect taxonomy covers six classes of dynamic visual anomaly
observable within a short interaction recording. Defects that manifest
only over longer sessions, that depend on backend state, or that are
purely functional rather than visual are out of scope.


\section{Dissertation Structure}
\label{sec:structure}

The remainder of this dissertation is organised as follows.
Chapter~\ref{ch:litreview} reviews work on visual regression testing, the
characterisation of web UI defects, the application of vision-language
models to visual software engineering tasks, and the construction of video
input for such models, concluding with an explicit statement of the
research gap. Chapter~\ref{chap:methodology} sets out the methodology,
covering dataset generation by automated fault injection, the keyframe
selection and dual-path decision components, the four-condition ablation
design, and the statistical protocol.
Chapter~\ref{chap:results_and_discussion} reports and discusses the
results. Chapter~\ref{chap:conclusion} concludes and identifies directions
for further work.